\documentclass{article}

\def\npart {IA}
\def\nterm {Michaelmas}
\def\nyear {2023}
\def\nlecturer {Prof N.\ Peake}
\def\ncourse {Vector and Matrices}

\usepackage{alltt}
\usepackage{amsmath, amsfonts, amssymb, amsthm}
\usepackage{booktabs}
\usepackage{caption}
\usepackage{enumitem}
\usepackage{fancyhdr}
\usepackage{graphicx}
\usepackage{mathdots}
\usepackage{mathtools}
\usepackage{microtype}
\usepackage{multirow}
\usepackage{pdflscape}
\usepackage{pgfplots}
\usepackage{siunitx}
\usepackage{textcomp}
\usepackage{slashed}
\usepackage{tabularx}
\usepackage{tikz}
\usepackage{tkz-euclide}
\usepackage[normalem]{ulem}
\usepackage[all]{xy}

\pgfplotsset{compat=1.18}
\pagestyle{fancyplain}

\lhead{\emph{\nouppercase{\leftmark}}}
\rhead{
  \ifnum\thepage=1
  \else
    \npart\ \vline\ \ncourse
  \fi}

\def\nauthor{Marcus Ng}
\author{Based on lectures by \nlecturer \\\small Notes taken by \nauthor}
\date{\nterm\ \nyear}
\title{Part \npart\ --- \ncourse}

\newcommand*{\Cdot}{{\raisebox{-0.25ex}{\scalebox{1.5}{$\cdot$}}}}
\setlist[enumerate,1]{label={(\roman*)}}

\newcommand {\pd}[2][ ]{
  \ifx #1 { }
    \frac{\partial}{\partial #2}
  \else
    \frac{\partial^{#1}}{\partial #2^{#1}}
  \fi
}

% Theorems
\theoremstyle{definition}
\newtheorem*{aim}{Aim}
\newtheorem*{axiom}{Axiom}
\newtheorem*{claim}{Claim}
\newtheorem*{cor}{Corollary}
\newtheorem*{conjecture}{Conjecture}
\newtheorem*{defi}{Definition}
\newtheorem*{eg}{Example}
\newtheorem*{ex}{Exercise}
\newtheorem*{fact}{Fact}
\newtheorem*{law}{Law}
\newtheorem*{lemma}{Lemma}
\newtheorem*{notation}{Notation}
\newtheorem*{prop}{Proposition}
\newtheorem*{question}{Question}
\newtheorem*{problem}{Problem}
\newtheorem*{rrule}{Rule}
\newtheorem*{thm}{Theorem}
\newtheorem*{assumption}{Assumption}
\newtheorem*{assert}{Assertion}

\newtheorem*{remark}{Remark}
\newtheorem*{warning}{Warning}
\newtheorem*{exercise}{Exercise}

\newtheorem{nthm}{Theorem}[section]
\newtheorem{nlemma}[nthm]{Lemma}
\newtheorem{nprop}[nthm]{Proposition}
\newtheorem{ncor}[nthm]{Corollary}

\SetEnumitemKey{cases}{
  label=\underline{Case~\arabic*:},
  ref={\arabic*},
  align=left,
  leftmargin=\parindent
}

% \renewcommand{\labelitemi}{--}
% \renewcommand{\labelitemii}{$\circ$}
% \renewcommand{\labelenumi}{(\roman{*})}

\let\stdsection\section
\renewcommand\section{\newpage\stdsection}

% Strike through
\def\st{\bgroup \ULdepth=-.55ex \ULset}


%%%%%%%%%%%%%%%%%%%%%%%%%
%%%%% Maths Symbols %%%%%
%%%%%%%%%%%%%%%%%%%%%%%%%

% Logical symbols
\newcommand{\contradiction}{
  \begin{tikzpicture}[rotate=45,x=0.5ex,y=0.5ex]
  \draw[line width=.1ex] (0,2) -- (3,2) (0,1) -- (3,1) (1,3) -- (1,0) (2,3) -- (2,0);
  \end{tikzpicture}
}


\let\U\relax
\let\C\relax
\let\G\relax

% Matrix groups
% \newcommand{\GL}{\mathrm{GL}}
% \newcommand{\Or}{\mathrm{O}}
% \newcommand{\PGL}{\mathrm{PGL}}
% \newcommand{\PSL}{\mathrm{PSL}}
% \newcommand{\PSO}{\mathrm{PSO}}
% \newcommand{\PSU}{\mathrm{PSU}}
% \newcommand{\SL}{\mathrm{SL}}
% \newcommand{\SO}{\mathrm{SO}}
% \newcommand{\Spin}{\mathrm{Spin}}
% \newcommand{\Sp}{\mathrm{Sp}}
% \newcommand{\SU}{\mathrm{SU}}
\newcommand{\U}{\mathrm{U}}
% \newcommand{\Mat}{\mathrm{Mat}}

% Matrix algebras
% \newcommand{\gl}{\mathfrak{gl}}
% \newcommand{\ort}{\mathfrak{o}}
% \newcommand{\so}{\mathfrak{so}}
% \newcommand{\su}{\mathfrak{su}}
% \newcommand{\uu}{\mathfrak{u}}
% \renewcommand{\sl}{\mathfrak{sl}}

% Special sets
\newcommand{\C}{\mathbb{C}}
\newcommand{\CP}{\mathbb{CP}}
\newcommand{\GG}{\mathbb{G}}
\newcommand{\N}{\mathbb{N}}
\newcommand{\Q}{\mathbb{Q}}
\newcommand{\R}{\mathbb{R}}
\newcommand{\RP}{\mathbb{RP}}
\newcommand{\T}{\mathbb{T}}
\newcommand{\Z}{\mathbb{Z}}
\renewcommand{\H}{\mathbb{H}}

% Brackets
\newcommand{\abs}[1]{\left\lvert #1\right\rvert}
\newcommand{\bket}[1]{\left\lvert #1\right\rangle}
\newcommand{\brak}[1]{\left\langle #1 \right\rvert}
\newcommand{\braket}[2]{\left\langle #1\middle\vert #2 \right\rangle}
\newcommand{\bra}{\langle}
\newcommand{\ket}{\rangle}
\newcommand{\norm}[1]{\left\lVert #1\right\rVert}
\newcommand{\normalorder}[1]{\mathop{:}\nolimits\!#1\!\mathop{:}\nolimits}
\newcommand{\tv}[1]{|#1|}
\renewcommand{\vec}[1]{\boldsymbol{\mathbf{#1}}}

% not-math
% \newcommand{\bolds}[1]{{\bfseries #1}}
% \newcommand{\cat}[1]{\mathsf{#1}}
% \newcommand{\ph}{\,\cdot\,}
% \newcommand{\term}[1]{\emph{#1}\index{#1}}
% \newcommand{\phantomeq}{\hphantom{{}={}}}

% Probability
% \DeclareMathOperator{\Bernoulli}{Bernoulli}
% \DeclareMathOperator{\betaD}{beta}
% \DeclareMathOperator{\bias}{bias}
% \DeclareMathOperator{\binomial}{binomial}
% \DeclareMathOperator{\corr}{corr}
% \DeclareMathOperator{\cov}{cov}
% \DeclareMathOperator{\gammaD}{gamma}
% \DeclareMathOperator{\mse}{mse}
% \DeclareMathOperator{\multinomial}{multinomial}
% \DeclareMathOperator{\Poisson}{Poisson}
% \DeclareMathOperator{\var}{var}
% \newcommand{\E}{\mathbb{E}}
% \newcommand{\Prob}{\mathbb{P}}

% Algebra
% \DeclareMathOperator{\adj}{adj}
% \DeclareMathOperator{\Ann}{Ann}
% \DeclareMathOperator{\Aut}{Aut}
% \DeclareMathOperator{\Char}{char}
% \DeclareMathOperator{\disc}{disc}
% \DeclareMathOperator{\dom}{dom}
% \DeclareMathOperator{\fix}{fix}
% \DeclareMathOperator{\Hom}{Hom}
% \DeclareMathOperator{\id}{id}
% \DeclareMathOperator{\image}{image}
% \DeclareMathOperator{\im}{im}
% \DeclareMathOperator{\re}{re}
% \DeclareMathOperator{\tr}{tr}
% \DeclareMathOperator{\Tr}{Tr}
% \newcommand{\Bilin}{\mathrm{Bilin}}
% \newcommand{\Frob}{\mathrm{Frob}}

% Others
% \newcommand\ad{\mathrm{ad}}
% \newcommand\Art{\mathrm{Art}}
% \newcommand{\B}{\mathcal{B}}
% \newcommand{\cU}{\mathcal{U}}
% \newcommand{\Der}{\mathrm{Der}}
% \newcommand{\D}{\mathrm{D}}
% \newcommand{\dR}{\mathrm{dR}}
% \newcommand{\exterior}{\mathchoice{{\textstyle\bigwedge}}{{\bigwedge}}{{\textstyle\wedge}}{{\scriptstyle\wedge}}}
% \newcommand{\F}{\mathbb{F}}
% \newcommand{\G}{\mathcal{G}}
% \newcommand{\Gr}{\mathrm{Gr}}
% \newcommand{\haut}{\mathrm{ht}}
% \newcommand{\Hol}{\mathrm{Hol}}
% \newcommand{\hol}{\mathfrak{hol}}
% \newcommand{\Id}{\mathrm{Id}}
% \newcommand{\lie}[1]{\mathfrak{#1}}
% \newcommand{\op}{\mathrm{op}}
% \newcommand{\Oc}{\mathcal{O}}
% \newcommand{\pr}{\mathrm{pr}}
% \newcommand{\Ps}{\mathcal{P}}
% \newcommand{\pt}{\mathrm{pt}}
% \newcommand{\qeq}{\mathrel{``{=}"}}
% \newcommand{\Rs}{\mathcal{R}}
% \newcommand{\Vect}{\mathrm{Vect}}
% \newcommand{\wsto}{\stackrel{\mathrm{w}^*}{\to}}
% \newcommand{\wt}{\mathrm{wt}}
% \newcommand{\wto}{\stackrel{\mathrm{w}}{\to}}
% \renewcommand{\d}{\mathrm{d}}
% \renewcommand{\P}{\mathbb{P}}
%\renewcommand{\F}{\mathcal{F}}

\let\Im\relax
\let\Re\relax

% \DeclareMathOperator{\area}{area}
% \DeclareMathOperator{\card}{card}
% \DeclareMathOperator{\ccl}{ccl}
% \DeclareMathOperator{\ch}{ch}
% \DeclareMathOperator{\cl}{cl}
% \DeclareMathOperator{\cls}{\overline{\mathrm{span}}}
% \DeclareMathOperator{\coker}{coker}
% \DeclareMathOperator{\conv}{conv}
% \DeclareMathOperator{\cosec}{cosec}
% \DeclareMathOperator{\cosech}{cosech}
% \DeclareMathOperator{\covol}{covol}
% \DeclareMathOperator{\diag}{diag}
% \DeclareMathOperator{\diam}{diam}
% \DeclareMathOperator{\Diff}{Diff}
% \DeclareMathOperator{\End}{End}
% \DeclareMathOperator{\energy}{energy}
% \DeclareMathOperator{\erfc}{erfc}
% \DeclareMathOperator{\erf}{erf}
% \DeclareMathOperator*{\esssup}{ess\,sup}
% \DeclareMathOperator{\ev}{ev}
% \DeclareMathOperator{\Ext}{Ext}
% \DeclareMathOperator{\fst}{fst}
% \DeclareMathOperator{\Fit}{Fit}
% \DeclareMathOperator{\Frac}{Frac}
% \DeclareMathOperator{\Gal}{Gal}
% \DeclareMathOperator{\gr}{gr}
% \DeclareMathOperator{\hcf}{hcf}
\DeclareMathOperator{\Im}{Im}
% \DeclareMathOperator{\Ind}{Ind}
% \DeclareMathOperator{\Int}{Int}
% \DeclareMathOperator{\Isom}{Isom}
% \DeclareMathOperator{\lcm}{lcm}
% \DeclareMathOperator{\length}{length}
% \DeclareMathOperator{\Lie}{Lie}
% \DeclareMathOperator{\like}{like}
% \DeclareMathOperator{\Lk}{Lk}
% \DeclareMathOperator{\Maps}{Maps}
% \DeclareMathOperator{\orb}{orb}
% \DeclareMathOperator{\ord}{ord}
% \DeclareMathOperator{\otp}{otp}
% \DeclareMathOperator{\poly}{poly}
% \DeclareMathOperator{\rank}{rank}
% \DeclareMathOperator{\rel}{rel}
% \DeclareMathOperator{\Rad}{Rad}
\DeclareMathOperator{\Re}{Re}
% \DeclareMathOperator*{\res}{res}
% \DeclareMathOperator{\Res}{Res}
% \DeclareMathOperator{\Ric}{Ric}
% \DeclareMathOperator{\rk}{rk}
% \DeclareMathOperator{\Rees}{Rees}
% \DeclareMathOperator{\Root}{Root}
% \DeclareMathOperator{\sech}{sech}
% \DeclareMathOperator{\sgn}{sgn}
% \DeclareMathOperator{\snd}{snd}
% \DeclareMathOperator{\Spec}{Spec}
% \DeclareMathOperator{\spn}{span}
% \DeclareMathOperator{\stab}{stab}
% \DeclareMathOperator{\St}{St}
% \DeclareMathOperator{\supp}{supp}
% \DeclareMathOperator{\Syl}{Syl}
% \DeclareMathOperator{\Sym}{Sym}
% \DeclareMathOperator{\vol}{vol}
\numberwithin{equation}{section}

\begin{document}
\maketitle{
    \small 
    \noindent\textbf{Complex numbers}\\
    Review of complex numbers, including complex conjugate, inverse, modulus, argument and Argand diagram. Informal treatment of complex logarithm, $n$-th roots and complex powers. de Moivre's theorem.\hspace*{\fill}[2]

    \vspace{10pt}
    \noindent\textbf{Vectors}\\
    Review of elementary algebra of vectors in $\R^3$, including scalar product. Brief discussion of vectors in $\R^n$ and $\C^n$; scalar product and the Cauchy-Schwarz inequality. Concepts of linear span, linear independence, subspaces, basis and dimension.

    \vspace{5pt}
    \noindent Suffix notation: including summation convention, $\delta_{ij}$ and $\varepsilon_{ijk}$. Vector product and triple product: definition and geometrical interpretation. Solution of linear vector equations. Applications of vectors to geometry, including equations of lines, planes and spheres.\hspace*{\fill}[5]

    \vspace{10pt}
    \noindent\textbf{Matrices}\\
    Elementary algebra of $3\times 3$ matrices, including determinants. Extension to $n\times n$ complex matrices. Trace, determinant, non-singular matrices and inverses. Matrices as linear transformations; examples of geometrical actions including rotations, reflections, dilations, shears; kernel and image.\hspace*{\fill}[4]

    \vspace{5pt}
    \noindent Simultaneous linear equations: matrix formulation; existence and uniqueness of solutions, geometric interpretation; Gaussian elimination.\hspace*{\fill}[3]

    \vspace{5pt}
    \noindent Symmetric, anti-symmetric, orthogonal, hermitian and unitary matrices. Decomposition of a general matrix into isotropic, symmetric trace-free and antisymmetric parts.\hspace*{\fill}[1]

    \vspace{10pt}
    \noindent\textbf{Eigenvalues and Eigenvectors}\\
    Eigenvalues and eigenvectors; geometric significance.\hspace*{\fill}[2]

    \vspace{5pt}
    \noindent Proof that eigenvalues of hermitian matrix are real, and that distinct eigenvalues give an orthogonal basis of eigenvectors. The effect of a general change of basis (similarity transformations). Diagonalization of general matrices: sufficient conditions; examples of matrices that cannot be diagonalized. Canonical forms for $2 \times 2$ matrices.\hspace*{\fill}[5]

    \vspace{5pt}
    \noindent Discussion of quadratic forms, including change of basis. Classification of conics, cartesian and polar forms.\hspace*{\fill}[1]

    \vspace{5pt}
    \noindent Rotation matrices and Lorentz transformations as transformation groups.\hspace*{\fill}[1]}

\tableofcontents

%%%%%%%%%%%%%%%%%%%%%%
% Lecture 1: 6/10/2023
%%%%%%%%%%%%%%%%%%%%%%
\section{Complex Numbers}
\subsection{Basic properties}
\begin{defi}[Complex number]
    A \emph{complex number} $z \in \C$ is
    \[
        z = a + ib, \quad a, b \in \R, \quad i = \sqrt{-1}  
    \]
    We define two other functions, $\Re()$ and $\Im()$ such that,
    \[
        \Re(z) = a \quad \Im(z) = b 
    \]
    Clearly, we may note that $\R \subset \C$ in the case where $b = 0$
\end{defi}

We can define algebraic manipulation over the complex numbers,
\begin{align*}
    z_1 \pm z_2 &= a_1 \pm a_2 + i(b_1 \pm b_2) \\
    z_1z_2 &= a_1a_2 - b_1b_2 + i(a_1b_2 + a_2b_1)
\end{align*}
We can define the inverse of a complex number
\begin{equation}
    z^{-1} = \frac{1}{a + ib} = \frac{a - ib}{a^2 + b^2}\label{eq:1}
\end{equation}
this demonstrates that $\C$ is \emph{closed} under addition and multiplication. We also define a complex conjugate of $z = a + ib$ as $\overline{z}$ or $z^*$.
\[
    \overline{z}, z^* = a - ib
\]
The modulus of $z$ as $\abs{z}$
\begin{align*}
    \abs{z} &= \abs{a + ib} = \sqrt{a^2 + b^2} \\
    z\overline{z} &= a^2 + b^2 = \abs{z}^2 \\
    z^{-1} &= \frac{\overline{z}}{\abs{z}^2} \tag{see \eqref{eq:1}}
\end{align*}
We also introduce the Argand Diagram, where the coordinate axes are the Real and Imaginary axes. We can see that $\overline{z}$ is the reflection of $z$ in the $x$-axis.

% draw an argand diagram in tikz
\begin{defi}[Triangle inequality]
    For $z_1, z_2 \in \C$, the triangle inequality states that
    \begin{equation}
        \abs{z_1 + z_2} \leq \abs{z_1} + \abs{z_2}\label{eq:2}
    \end{equation}
\end{defi}

\begin{prop}
    The triangle inequality has another common alternative form,
    \begin{equation}
        \abs{z_1 - z_2} \geq \abs{\abs{z_1} - \abs{z_2}}\label{eq:3}
    \end{equation}
\end{prop}

\begin{proof}
    Write $z_1 = z^{\prime}_1 - z^{\prime}_2$, and $z_2 = z^{\prime}_2$ and substitute them into \eqref{eq:2}. Then,
    \begin{align*}
        &\abs{z^{\prime}_1} \leq \abs{z^{\prime}_1 - z^{\prime}_2} + \abs{z^{\prime}_2} \\
        \Rightarrow& \abs{z^{\prime}_1} - \abs{z^{\prime}_2} \leq\abs{z^{\prime}_1 - z^{\prime}_2}\\
        \Rightarrow& \abs{z^{\prime}_1 - z^{\prime}_2} \geq \abs{z^{\prime}_1} - \abs{z^{\prime}_2} \tag{$\star$}
    \end{align*}
    Notice that our labelling of $z_1$ and $z_2$ was arbitrary, and so we can freely swap the assignment. This gives,
    \[
        \abs{z^{\prime}_1 - z^{\prime}_2} \geq \abs{z^{\prime}_2} - \abs{z^{\prime}_1} \tag{$\star\star$}
    \]
    Since $\abs{z^{\prime}_1 - z^{\prime}_2} \geq (\star) \ \text{and} \ (\star\star)$, hence, \eqref{eq:3}
\end{proof}

\begin{defi}[Polar Representation]
    For a complex number $z = a + ib$ we use the argand diagram to write $z = r(\cos{\theta} + i\sin{\theta})$, where
    \[
        r = \abs{z} \quad \text{and} \quad \theta = \arctan\left(\frac{b}{a}\right) = \arg{(z)}  
    \]
\end{defi}

\begin{remark}
    Note that the mapping $(r, \theta) \rightarrow z$ is unique. However, the reverse $r \rightarrow (r, \theta)$ is \emph{NOT} unique since for each value of $\theta$, $\theta + 2n\pi$ is equally valid.

    In order to make this mapping single valued, we will by convention restrict $\theta$ to the principle value, where
    \[
        -\pi < \theta \leq \pi  
    \]

    This multivalued property becomes quite interesting when we consider complex functions. We can imagine taking the multivalued imaginary $\ln$ function as a multistory carpark, walking round and round up and down. To fix this what we do is to build a brick wall on the upramp (``branch cut'').
\end{remark}

Using complex numbers in this form allows us to view complex arithmetic geometrically. For instance, let
\begin{align*}
    z_k &= x_k + iy_k \quad \text{for} \quad k = 1, 2 \\
    &= r_k(\cos{\theta_k} + i\sin(\theta_k))
\end{align*}
\begin{align*}
    z_1z_2 &= r_1(\cos{\theta_1} + i\sin{\theta_1})r_2(\cos{\theta_2} + i\sin{\theta_2}) \\
    &= r_1r_2[(\cos{\theta_1}\cos{\theta_2} - \sin{\theta_1}\sin{\theta_2}) + i(\cos{\theta_1}\sin{\theta_2} + \sin{\theta_1}\cos{\theta_2})] \\
    &= r_1r_2(\cos{(\theta_1 + \theta_2)} + i\sin{(\theta_1 + \theta_2)}) \numberthis \label{eq:4}
\end{align*}

\subsection{Complex exponential function}
\begin{defi}[Complex exponential function]
    \begin{align*}
        \exp(z) &= 1 + z + \frac{z^2}{2!} + \cdots + \frac{z^n}{n!} + \cdots \\
        &=\sum_{n=0}^{\infty}{\frac{z^n}{n!}} \numberthis \label{eq:5}
    \end{align*}
\end{defi}
This function is ``uniformly convergent'' (something that we will cover next term in analysis). That means that we are able to do lots of algebraic manipulation without too much hesitation. =

\begin{prop}[Multiplication of $\exp()$]
    \[
        \exp(z_1)\exp(z_2) = \exp(z_1 + z_2)  
    \]
\end{prop}
\begin{proof}
    We will consider the double sum
    \begin{align*}
        \sum_{m,n=0}^{\infty}{a_mn} &= a_{00} + a_{01} + a_{02} + a_{03} + \cdots \\
         &+ a_{10} + a_{11} + a_{12} + \cdots \\
         &+ a_{20} + a_{21} + a_{22} + \cdots \\
         &+ a_{30} + a_{31} + a_{32} + \cdots \\
         &\vdots
    \end{align*}
    Then by collecting terms along the diagonal we can regroup the sum into
    \[
        =\sum_{r=0}^{\infty}\sum_{m=0}^{r}{a_{r-m, m}}
    \]
    Hence, 
    \begin{align*}
        \exp(z_1) \cdot \exp(z_2) &= \sum_{r=0}^{\infty}\sum_{m=0}^{r}{\frac{z_1^{r-m}}{(r-m)!}\frac{z_2^{m}}{m!}} \\
        &= \sum_{r=0}^{\infty}{\frac{1}{r!}}\sum_{m=0}^{r}{\frac{r!}{(r-m)!m!}z_1^{r-m}z_2^{m}} \\
        &= \sum_{r=0}^{\infty}{\frac{1}{r!} (z_1 + z_2)^r} \\
        &= \exp(z_1 + z_2)
    \end{align*}
\end{proof}

%%%%%%%%%%%%%%%%%%%%%%
% Lecture 2: 8/10/2023
%%%%%%%%%%%%%%%%%%%%%%
Similarly, we can define the complex $\sin(z)$ and $\cos(z)$ functions.
\begin{defi}[Complex sine and cosine]
    \begin{align*}
        \sin{z} &= \sum_{n=0}^{\infty}{(-1)^n\frac{z^{2n+1}}{(2n+1)!}} \\
        \cos{z} &= \sum_{n=0}^{\infty}{(-1)^n\frac{z^{2n}}{(2n)!}} \numberthis \label{eq:6}
    \end{align*}
\end{defi}

From \eqref{eq:5} we can derive that
\begin{align*}
    e^{iz} &= \sum_{n=0}^{\infty}{\frac{i^nz^n}{n!}} \\
    &= \sum_{n=0}^{\infty}{(-1)^n\frac{z^{2n+1}i^{2n+1}}{(2n+1)!}} + \sum_{n=0}^{\infty}{(-1)^n\frac{z^{2n}i^{2n}}{(2n)!}} \\
    &= i\sin{z} + \cos{z} \numberthis \label{eq:7}
\end{align*}

We may consider $\theta \in \R$ which gives,
\begin{equation}
    z = re^{i\theta} = r(\cos{\theta} + i\sin{\theta}) \label{eq:8}
\end{equation}

\subsection{Roots of unity}
We seek to solve equations of the form $z^n = 1$ where $n \in \Z$. Since this is a polynomial of degree $n$ we expect to find $n$ solutions. By applying the exponential form we found earlier,
\begin{align*}
    z^n = 1 &= \exp(i2 \pi k), \quad k \in \Z \\
    z &= \exp(i2 \pi \frac{k}{n})
\end{align*}
Notice that the roots will repeat after $k = n-1$. Hence the solutions are
\[
    z = \exp(i2 \pi \frac{k}{n}), \quad k = 0, 1, 2, \cdots, n-1
\]

\begin{eg}[$ z^7 = 2$]
    \[
        z = 2^{\frac{1}{7}} \exp(i 2\pi \frac{m}{7}), \quad m = 0, 1, 2, \cdots, 6
    \]
\end{eg}

An interesting point to note is that for the polynomial $z^n - 1 = 0$, if we arbitrarily choose a root $w = \exp(\frac{2\pi i}{n})$, then all the remaining roots will be
\[
    1, w, w^2, w^3, \cdots, w^{n-1}
\]
Further, if we note that since the coefficient of $z^{n-1}$ is $0$, that implies that the sum of the roots must be $0$.
\[
   1 + w + w^2 + w^3 + \cdots + w^{n-1} = 0
\]

\subsection{Complex Logarithm and Powers}
\begin{defi}[Complex Log]
    The \emph{complex log}, $\log{z}$, is a solution of the following
    \[
        e^w = z  
    \]
    Hence, $\exp()$ and $\log()$ are inverse functions, i.e;
    \[
        e^{\log{z}} = z
    \]
    If we write $z = re^{i\theta}$ from \eqref{eq:8} then,
    \begin{align*}
        \log{z} = \log{re^{i\theta}} &= \log{r} + \log{e^{i\theta}} \\
        &= log{\abs{z}} + i\arg{(z)} \numberthis \label{eq:9}
    \end{align*}
    It is important to note that this is not a single valued function. In order to define a single valued $\log$ we will insist that
    \[
        -\pi < \arg{(z)} \leq \pi
    \]
    Hence,
    \begin{equation}
        \log{z} = log{\abs{z}} + i\theta, \quad -\pi < \theta \leq \pi \label{eq:10}
    \end{equation}
\end{defi}

\begin{eg}
    \begin{align*}
        \log{(2i)} &= \log{2e^{i\frac{\pi}{2}}} \\
        &=\log{2} + i\frac{\pi}{2}
    \end{align*}
\end{eg}

\begin{defi}[Complex Power]
    Since we have defined the complex logarithm we can now define the complex power, $z^\alpha$ for $z, \alpha \in \C$ as
    \begin{equation}
        z^\alpha = e^{\alpha \log{z}} \label{eq:11}
    \end{equation}
    Again, this function is multivalued. Specifically using \eqref{eq:10},
    \[
        z^\alpha = e^{\alpha \log{z}}e^{i \alpha \theta}e^{i \alpha 2\pi n}, \quad n \in \Z
    \]
    However, unlike the roots of unity, there is no reason why this series will repeat and terminate. If $\alpha$ is real and rational then will there be finite solutions. Once again, we will make this function single valued by insisting on using the principle value for $\theta$
    \[
        -\pi < \arg{(z)} \leq \pi
    \]
\end{defi}

\begin{ex}\leavevmode
    \begin{enumerate}
        \item \begin{align*}
            (1 + i)^\frac{1}{2} &= e^{\frac{1}{2}\sqrt{2}}e^\frac{i\pi}{8}e^{in\pi}, \quad n = 0, 1 \\
            &= 2^\frac{1}{4}e^\frac{i\pi}{8}, 2^\frac{1}{4}e^\frac{i9\pi}{8} \\
            &= 2^\frac{1}{4}e^\frac{i\pi}{8} \tag{by insiting on the principle value}
        \end{align*}
        \item \begin{align*}
            i^i &= \exp(i \log{(e^\frac{i\pi}{2})}) \tag{from \eqref{eq:11}} \\
            &= \exp(i \frac{i\pi}{2}) \\
            &= \exp(-\frac{\pi}{2})
        \end{align*}
    \end{enumerate}
\end{ex}

\subsection{De Moivre's Theorem}
\begin{thm}[De Moivre's Theorem]
    \begin{equation}
        \cos{n\theta} + i\sin{n\theta} = (\cos{\theta} + i\sin{\theta})^n, \quad n \in \Z \label{eq:12}
    \end{equation}
\end{thm}

\begin{proof}
    By induction, we will first consider the case where $n > 0$. Clearly, $n=1$ is true. Assume $n=k$ is true for $k > 1$.
    \begin{align*}
        (\cos{\theta} + i\sin{\theta})^{k+1} &= (\cos{k\theta} + i\sin{k\theta})(\cos{\theta} + i\sin{\theta}) \\
        &= \cos{k\theta}\cos{\theta} - \sin{k\theta}\sin{\theta} + i(\cos{k\theta}\sin{\theta} + \sin{k\theta}\cos{\theta}) \\
        &= \cos{(k+1)\theta} + i\sin{(k+1)\theta}
    \end{align*}
    Now when $n < 0$, we write $m = -n$ such that $m > 0$
    \begin{align*}
        (\cos{\theta} + i\sin{\theta})^{-m} &= (\cos{m\theta} + i\sin{m\theta})^{-1} \\
        &= \cos{m\theta} - i\sin{m\theta} \tag{by ($\star$)}\\
        &= \cos{(-m\theta)} - i\sin{(-m\theta)}
    \end{align*}
    Note $(\star)$:
    \[
        z^{-1} = \frac{z^*}{\abs{z}^2}    
    \]
\end{proof}

Using De Moivre's theorem we can work out multiple angle formulae by comparing real and imaginary parts:
\begin{eg}
    \begin{align*}
        \cos{5 \theta} + i\sin{5 \theta} &= (\cos{\theta} + i\sin{\theta})^5 \\
        &\vdots \\
        \cos{5 \theta} &= 5\cos{\theta} - 20 \cos^3{\theta} + 16\cos^5{\theta}\\
        \sin{5 \theta} &= 5\cos{\theta} - 20 \cos^3{\theta} + 16\cos^5{\theta}\\
    \end{align*}
\end{eg}
Note that for non-integer powers of $n$ we will have to insist on the principle value \eqref{eq:12}

\subsection{Lines and Circles in $\C$}
\begin{defi}[Straight Line]
    A line $l$ through $z_0 \in \C$ parallel to $w \in \C$ can be generally written
    \[
        z = z_0 + \lambda w, \quad \lambda \in \R    
    \]
    However, it would be better to get rid of the parameterisation. We can use the fact that $\lambda$ must be real and hence its conjugate is equal to itself.
    \begin{align*}
        \lambda &= \frac{z - z_0}{w} \in \R\\
        &\Rightarrow \frac{z - z_0}{w} = \frac{\overline{z} - \overline{z_0}}{\overline{w}} \\
        &\Leftrightarrow z \overline{w} - \overline{z}w = z_0 \overline{w} - \overline{z_0} w
    \end{align*}
\end{defi}

\begin{defi}[Circle]
    A circle can be defined from its center $c \in \C$ and radius $\rho \in \R$. By definition, if $z$ lies on the circle,
    \begin{align*}
        &\abs{z - c} = \rho \\
        \Rightarrow& (z - c)(\overline{z} - \overline{c}) = \rho^2 \\
        \Leftrightarrow& z\overline{z} - \overline{c}z - c\overline{z} = \rho^2 - c\overline{c}
    \end{align*}
\end{defi}

%%%%%%%%%%%%%%%%%%%%%%
% Lecture 3: 11/10/2023
%%%%%%%%%%%%%%%%%%%%%%

\section{Vectors}
\subsection{Definition \& basic properties}
\begin{defi}[Vector]
    A vector $\vec{x}$ has a positive length and a direction. If,
    \[
        \abs{\vec{x}} = 0 \Rightarrow \vec{x} = \vec{0}  
    \]
    If $\abs{\vec{x}} = 1$ then we write $\uvec{x}$ as \emph{unit vector}. Note the distinction between scalar $0$ and vector $\vec{0}$.
\end{defi}

\begin{defi}[Vector Field]
    A vector field $F = F(\vec{x})$ is a vector function of position and maybe time.
\end{defi}

\begin{defi}[Vector addition]
    % insert paralellogram diagram.
    We can define the addition of vectors using the paralellogram rule,
    \begin{align*}
        \vec{a} + \vec{b} &= \vec{c} \\
        \pvec{OA} + \pvec{OB} &= \pvec{OC}
    \end{align*}
    More abstractly, addition will have the following properties
    \begin{enumerate}[eqn]
        \item $\vec{a} + \vec{b} = \vec{b} + \vec{a}$ \hspace*{\fill}(commutativity) \label{eq:2-1}
        \item $\vec{a} + (\vec{b} + \vec{c}) = (\vec{a} + \vec{b}) + \vec{c}$ \hspace*{\fill}(associativity) \label{eq:2-2}
        \item $\exists \vec{0}, \quad \vec{a} + \vec{0} = \vec{a}$ \hspace*{\fill}(unique zero) \label{eq:2-3}
        \item $\vec{a} + -\vec{a} = 0$ \hspace*{\fill}(unique additive inverse) \label{eq:2-4}
    \end{enumerate}
\end{defi}
\begin{defi}[Multiplication by a scalar]
    For $\lambda \in \R$ (we will do complex scalars later in the course), $\lambda\vec{a}$ has magnitutde $\abs{\lambda}\abs{\vec{a}}$ and is parallel/anti-parallel to $\vec{a}$ if $\lambda<0$ \textbackslash $\lambda>0$. 
    Likewise, we have the following properties
    \begin{enumerate}[eqn]
        \item $(\lambda + \mu)\vec{a} = \lambda\vec{a} + \mu\vec{a}$ \label{eq:2-5}
        \item $\lambda(\vec{a} + \vec{b}) = \lambda\vec{a} + \lambda\vec{b}$ \label{eq:2-6}
        \item $\lambda(\mu\vec{a}) = (\lambda\mu)\vec{a}$ \label{eq:2-7}
        \item $1\vec{a} = \vec{a}$ \label{eq:2-8}
    \end{enumerate}
\end{defi}

Implicit in the two definitions is that the set of vectors is closed under addition and scalar multiplication.

\begin{defi}[Vector Space]
    $V$ is a \emph{vector space} if it is a set of vectors which satisfy properties \eqref{eq:2-1} -- \eqref{eq:2-8} inclusive and is closed under addition and multiplication. Note that for now, we will be considering the vector space $V$ over the real numbers $\R$, but we will consider complex fields later. 
\end{defi}

\begin{eg}[$\R^n$ is a vector space]
\end{eg}
\begin{proof}
    \[
        \R^n = \{\vec{x} = (x_1, x_2, \ldots, x_n) \mid x_j \in \R, \ 1 \leq j \leq n\}
    \]
    Consider $\vec{x} , \vec{y} \in \R$
    \[
        \vec{x} + \vec{y} = (x_1 + y_1, x_2 + y_2, \ldots, x_n + y_n,)
    \]
    From this we see that property \eqref{eq:2-1} and \eqref{eq:2-2} is satisfied as we can inherit them from the Real numbers.

    \begin{enumerate}
        \item $\vec{0} = (0, 0, \ldots, 0)$ \hspace*{\fill}\eqref{eq:2-3}
        \item $-\vec{x} = (-x_1, -x_2, \ldots, -x_n)$ \hspace*{\fill}\eqref{eq:2-4}
        \item $\lambda\vec{x} = (\lambda x_1, \lambda x_2, \ldots, \lambda x_n)$ \hspace*{\fill}\eqref{eq:2-5} - \eqref{eq:2-8}
    \end{enumerate}
\end{proof}

Geometrically, we may interprit $\R^1$ as a line, and $\R^2$ as a plane. However, notice that the converse is not always true. Due to \eqref{eq:2-3} the space has to contain $0$. Therefore, if the line or plane does not go through the origin, it is not a vector space.

\subsection{Scalar Product}
\begin{defi}[Geometric definition in $\R^2$, $\R^3$]
    % insert diagram of dot product
    Given two vectors $\vec{a}, \vec{b}$ we define
    \[
        \vec{a} \cdot \vec{b} = \abs{\vec{a}}\abs{\vec{b}}\cos{\theta}  
    \]
    where $\theta$ is the angle between the two vectors.
\end{defi}
From this, we can see that
\begin{enumerate}
    \item $\vec{a} \cdot \vec{b} = \vec{b} \cdot \vec{a}$ 
    \item $\vec{a} \cdot \vec{a} = 0 \Leftrightarrow \vec{a} = \vec{0}$
    \item $\vec{a} \cdot \vec{b} = 0 \Leftrightarrow$ $\vec{a}$ and $\vec{b}$ are orthogonal
    \[\numberthis\]
\end{enumerate}

\begin{defi}[Projection]
    % insert diagram of projection
    $\vec{b}^\perp$ is the projection of $\vec{b}$ onto $\vec{a}$. Specifically, 
    \begin{equation}
        \vec{b}^\perp = \abs{\vec{b}}\cos{\theta}\frac{\vec{a}}{\abs{\vec{a}}} = \frac{(\vec{b}\cdot\vec{a})\vec{a}}{(\vec{a}\cdot\vec{a})} = (\uvec{a} \cdot \vec{b})\uvec{a} 
    \end{equation}
\end{defi}

We will move to a most abstract defintion of a scalar product
\begin{defi}[Inner product]
    We define $\vec{x} \cdot \vec{y} = \braket{\vec{x}}{\vec{y}} \in \R$ as an \emph{inner product}, if it satisfies
    \begin{enumerate}[eqn]
        \item $\vec{x} \cdot \vec{y} = \vec{y} \cdot \vec{x}$ \hspace*{\fill}(Symmetry) \label{eq:2-11}
        \item $\vec{x}, \vec{y}, \vec{z} \in V, \quad, \lambda, \mu \in \R$ \hspace*{\fill}(Linearity in the second argument) \label{eq:2-12}
        \[
            \vec{x} \cdot (\lambda\vec{y} + \mu\vec{z}) = \lambda\vec{x}\vec{y} + \mu\vec{x}\vec{z}  
        \]
        \item $\vec{x} \cdot \vec{x} \geq 0 $ and $\vec{x} \cdot \vec{x} = 0 \Leftrightarrow \vec{x} = \vec{0}$ \hspace*{\fill}(Positive definite) \label{eq:2-13}
    \end{enumerate}

    Note that we can use \eqref{eq:2-11} together with \eqref{eq:2-12} to derive linearlity in the first argument,
    \begin{equation}
        (\lambda\vec{y} + \mu\vec{z}) \cdot \vec{x} = \lambda\vec{x}\vec{y} + \mu\vec{x}\vec{z} \label{eq:2-15}
    \end{equation}
    We will see later in the course that linearity in the first argument does not hold for vector spaces over the complex field.

    We can also define the \emph{norm} of a vector $\vec{a}$ as 
    \[
        (\vec{a} \cdot \vec{a})^\frac{1}{2} = \abs{\vec{a}} = \norm{\vec{a}}  
    \]
\end{defi}

\begin{remark}
    Using this abstracted definition of an inner product, we can define an inner product between functions. For example define for function $f$ and $g$
    \[
        \braket{f}{g} = \int_{0}^{1}{fg \ \diff x}
    \]
    This definition satisfies \eqref{eq:2-11} and \eqref{eq:2-12} and we can check \eqref{eq:2-13}
    \[
        \braket{f}{f} = \int_{0}^{1}{f^2 \ \diff x} \geq 0
    \]
    and further if 
    \[
        \int_{0}^{1}{f^2 \ \diff x} = 0 \Rightarrow f = 0
    \]
\end{remark}

\subsection{Cauchy-Schwartz Inequality}
\begin{thm}[CS Inequality]
    Given two vectors $\vec{x}$, $\vec{y}$
    \begin{equation}
        \vec{x} \cdot \vec{y} \leq \abs{\vec{x}}\abs{\vec{y}} \quad \vec{x}, \vec{y} \in \R^n
    \end{equation}
    with equality if and only if $\vec{x} = 0$ or $\vec{y} = 0$ or $\vec{x} = \lambda \vec{y}$ for some $\lambda \in \R$.
\end{thm}
\begin{proof}
    Consider $\abs{\vec{x} - \lambda \vec{y}}^2 \geq 0$
    \begin{align*}
        &(\vec{x} - \lambda \vec{y}) \cdot (\vec{x} - \lambda \vec{y}) \geq 0 \\
        \Rightarrow& \lambda^2 \vec{y} \cdot \vec{y} - 2\lambda\vec{x}\vec{y} + \vec{x} \cdot \vec{x} \geq 0
    \end{align*}
    Suppose $\vec{x}, \vec{y} \neq \vec{0}$, we have a quadratic equation in $\lambda$. Therefore since the quadratic is greater than or equal to $0$, the discriminant must be less than or equal to $0$.
    \begin{align*}
        \Delta &= [2(\vec{x} \cdot \vec{y})]^2 - 4 (\vec{y} \cdot \vec{y})(\vec{x} \cdot \vec{x}) \leq 0 \\
        \Rightarrow& [2(\vec{x} \cdot \vec{y})]^2 \leq 4 (\vec{y} \cdot \vec{y})(\vec{x} \cdot \vec{x}) \\
        \Rightarrow& (\vec{x} \cdot \vec{y})^2 \leq (\vec{y} \cdot \vec{y})(\vec{x} \cdot \vec{x}) \\
        \Rightarrow& (\vec{x} \cdot \vec{y})^2 \leq \abs{\vec{x}}^2\abs{\vec{y}}^2
    \end{align*}
    Furthermore, we have equality when
    \[
        \abs{\vec{x} - \lambda \vec{y}}^2 = 0 \Leftrightarrow \vec{x} = \lambda \vec{y}
    \]
    Hence the result
\end{proof}

\begin{eg} $\vec{x} = (\alpha, \beta, \gamma), y = (1, 1, 1)$
    \begin{align*}
        \alpha + \beta + \gamma &\leq \sqrt{\alpha^2 + \beta^2 + \gamma^2} \sqrt{3} \tag{\ref{eq:2-15}} \\
        \alpha^2 + \beta^2 + \gamma^2 &\geq \alpha\beta + \beta\gamma + \gamma\alpha \tag{rearranging}
    \end{align*}
\end{eg}

Note that we must not use the geometric definition of the dot product to prove the CS inequality because $\cos{\theta}$ is not defined for higher dimensions. However, we can use CS to define an angle in $\R^n$
\[
    \theta = \arccos{\left(\frac{\vec{x}\vec{y}}{\abs{\vec{x}}\abs{\vec{y}}}\right)}    
\]
since CS gives us that the fraction is less than or equal to $1$.
\end{document}